% GENERATED FILE. Edit canonical lessons, then run npm run generate:books.
% canonical-source-hash: fnv1a64:4b9ca2b8cc0dd5a5
% canonical-lessons: ZH-C366-bimian, ZH-C366-dangxin, ZH-C366-jianchi, ZH-C366-fangqi, ZH-C366-zhide

\chapter{Take My Advice}
\label{ch:zh-take-my-advice}
\hlchaptermodality{366}

\begin{chapteropening}
\textbf{By the end of this chapter:} \emph{I can tell someone to watch out, keep going and not give up, and say something is worth it.}

Complete the last lesson of chapter 366: I can tell someone to watch out, keep going and not give up, and say something is worth it.
\end{chapteropening}

\section[bìmiǎn]{bìmiǎn — to avoid}

\label{lesson:ZH-C366-bimian}



\begin{quote}
Before the new one: say the Mandarin for might as well, then the Mandarin for a word of advice.
\end{quote}

\subsection*{You'll want to know: bìmiǎn}
\emph{bìmiǎn} — \textquotedblleft{}to avoid\textquotedblright{}.

\begin{quote}
\emph{wǒmen yīnggāi bìmiǎn chūcuò} — We should try not to slip up.
\end{quote}

Its characters come later; for now it is a word for the ear.

\begin{cousinweb}
\emph{bìmiǎn} is built from \emph{bì} and \emph{miǎn}, both inherited from Middle Chinese, \emph{bjieH} and \emph{mjenX} (Guangyun, in Baxter's transcription), which glosses the first 'to turn away' and the second 'to stop' and 'to escape'.
\end{cousinweb}

\begin{grammarlens}[title={What you've built}]
One more word for giving reasons.
\end{grammarlens}

\subsection*{Your turn}
\begin{itemize}
\raggedright
  \item \emph{Say it:} \emph{bìmiǎn}
  \item \emph{Say it:} \emph{bìmiǎn}, once more
  \item \emph{Say it:} \emph{wǒmen yīnggāi bìmiǎn chūcuò} — We should try not to slip up.
  \item \emph{Recall:} say the Mandarin for might as well, then the Mandarin for a word of advice, then say \emph{bìmiǎn} again
\end{itemize}

\begin{tcolorbox}[breakable,colback=teal!4,colframe=teal!35!black,title={Before you move on}]
What does bìmiǎn mean? (\textquotedblleft{}To avoid\textquotedblright{}.) Say it once more.
\end{tcolorbox}
\section[\texorpdfstring{\zh{当心} (dāngxīn)}{当心 (dāngxīn)}]{\zh{当心} (dāngxīn) — to watch out; to look out}

\label{lesson:ZH-C366-dangxin}



\begin{quote}
Before the new one: say the Mandarin for a word of advice, then the Mandarin for to avoid.
\end{quote}

\subsection*{You'll want to know: \zh{当心}}
\textbf{\zh{当心}} — \emph{dāngxīn} — \textquotedblleft{}to watch out; to look out\textquotedblright{}.

\begin{quote}
\textbf{\zh{下雨的时候开车必须当心}} \emph{xiàyǔ de shíhou kāichē bìxū dāngxīn} — You must take care driving in the rain.
\end{quote}

\begin{cousinweb}
\textbf{\zh{当心}} (\emph{dāngxīn}) is built from \zh{当} and \zh{心} 'the heart', both inherited from Middle Chinese, \emph{tang} and \emph{sim} (Guangyun, in Baxter's transcription), which glosses the first 'to match' and 'to preside'. How the pair came to mean watching out is not settled.
\end{cousinweb}

\begin{grammarlens}[title={What you've built}]
One more word for giving reasons.
\end{grammarlens}

\subsection*{Your turn}
\begin{itemize}
\raggedright
  \item \emph{Say it:} \emph{dāngxīn}
  \item \emph{Say it:} \emph{dāngxīn}, once more
  \item \emph{Say it:} \emph{xiàyǔ de shíhou kāichē bìxū dāngxīn} — You must take care driving in the rain.
  \item \emph{Recall:} say the Mandarin for a word of advice, then the Mandarin for to avoid, then say \emph{dāngxīn} again
\end{itemize}

\begin{tcolorbox}[breakable,colback=teal!4,colframe=teal!35!black,title={Before you move on}]
What does \zh{当心} mean? (\textquotedblleft{}To watch out; to look out\textquotedblright{}.) Say it once more.
\end{tcolorbox}
\section[jiānchí]{jiānchí — to keep at it; to persist}

\label{lesson:ZH-C366-jianchi}



\begin{quote}
Before the new one: say the Mandarin for to avoid, then the Mandarin for to watch out.
\end{quote}

\subsection*{You'll want to know: jiānchí}
\emph{jiānchí} — \textquotedblleft{}to keep at it; to persist\textquotedblright{}.

\begin{quote}
\emph{wǒ juéde xué Hànyǔ yào jiānchí} — I think learning Chinese takes perseverance.
\end{quote}

Its characters come later; for now it is a word for the ear.

\begin{cousinweb}
\emph{jiānchí} is built from \emph{jiān} and \emph{chí}, both inherited from Middle Chinese, \emph{ken} and \emph{dri} (Guangyun, in Baxter's transcription), which glosses the first 'firm' and the second 'to hold'.
\end{cousinweb}

\begin{grammarlens}[title={What you've built}]
One more word for giving reasons.
\end{grammarlens}

\subsection*{Your turn}
\begin{itemize}
\raggedright
  \item \emph{Say it:} \emph{jiānchí}
  \item \emph{Say it:} \emph{jiānchí}, once more
  \item \emph{Say it:} \emph{wǒ juéde xué Hànyǔ yào jiānchí} — I think learning Chinese takes perseverance.
  \item \emph{Recall:} say the Mandarin for to avoid, then the Mandarin for to watch out, then say \emph{jiānchí} again
\end{itemize}

\begin{tcolorbox}[breakable,colback=teal!4,colframe=teal!35!black,title={Before you move on}]
What does jiānchí mean? (\textquotedblleft{}To keep at it; to persist\textquotedblright{}.) Say it once more.
\end{tcolorbox}
\section[fàngqì]{fàngqì — to give up}

\label{lesson:ZH-C366-fangqi}



\begin{quote}
Before the new one: say the Mandarin for to watch out, then the Mandarin for to keep at it.
\end{quote}

\subsection*{You'll want to know: fàngqì}
\emph{fàngqì} — \textquotedblleft{}to give up\textquotedblright{}.

\begin{quote}
\emph{nǐ bù yīnggāi fàngqì} — You shouldn't give up.
\end{quote}

Its characters come later; for now it is a word for the ear.

\begin{cousinweb}
\emph{fàngqì} is built from \emph{fàng} and \emph{qì}, both inherited from Middle Chinese, \emph{pjangH} and \emph{khjijH} (Guangyun, in Baxter's transcription), which glosses the first 'to drive out' and the second, quoting an older dictionary, 'to throw away'.
\end{cousinweb}

\begin{grammarlens}[title={What you've built}]
One more word for giving reasons.
\end{grammarlens}

\subsection*{Your turn}
\begin{itemize}
\raggedright
  \item \emph{Say it:} \emph{fàngqì}
  \item \emph{Say it:} \emph{fàngqì}, once more
  \item \emph{Say it:} \emph{nǐ bù yīnggāi fàngqì} — You shouldn't give up.
  \item \emph{Recall:} say the Mandarin for to watch out, then the Mandarin for to keep at it, then say \emph{fàngqì} again
\end{itemize}

\begin{tcolorbox}[breakable,colback=teal!4,colframe=teal!35!black,title={Before you move on}]
What does fàngqì mean? (\textquotedblleft{}To give up\textquotedblright{}.) Say it once more.
\end{tcolorbox}
\section[zhíde]{zhíde — to be worth it; to be worth doing}

\label{lesson:ZH-C366-zhide}



\begin{quote}
Before the new one: say the Mandarin for to keep at it, then the Mandarin for to give up.
\end{quote}

\subsection*{You'll want to know: zhíde}
\emph{zhíde} — \textquotedblleft{}to be worth it; to be worth doing\textquotedblright{}.

\begin{quote}
\emph{zhè běn shū zhíde kàn} — This book is worth reading.
\end{quote}

Its characters come later; for now it is a word for the ear.

\begin{cousinweb}
\emph{zhíde} is built from \emph{zhí} and \zh{得} 'to get', both inherited from Middle Chinese, \emph{driH} and \emph{tok} (Guangyun, in Baxter's transcription), which glosses the first 'to hold' and 'to meet with', and the second 'gain'. The second is said light here.
\end{cousinweb}

\begin{grammarlens}[title={What you've built}]
One more word for giving reasons.
\end{grammarlens}

\subsection*{Your turn}
\begin{itemize}
\raggedright
  \item \emph{Say it:} \emph{zhíde}
  \item \emph{Say it:} \emph{zhíde}, once more
  \item \emph{Say it:} \emph{zhè běn shū zhíde kàn} — This book is worth reading.
  \item \emph{Recall:} say the Mandarin for to keep at it, then the Mandarin for to give up, then say \emph{zhíde} again
\end{itemize}

\begin{tcolorbox}[breakable,colback=teal!4,colframe=teal!35!black,title={Before you move on}]
What does zhíde mean? (\textquotedblleft{}To be worth it; to be worth doing\textquotedblright{}.) Say it once more.
\end{tcolorbox}
